% ============================================================
% Data Sources and Models — proposal_18E/data_sources.tex
% Reviewer concern: RF-5b (Data Sources TODO)
% Verified by: AC-6
%
% This file currently reflects PLAN B (public corpora) as a safe default.
% PLAN A (specific Sandia Kokkos/Trilinos/LAMMPS corpus names) will be
% swapped in after the 2026-04-09 Sandia meeting — see \todo{} marker below.
% ============================================================

% [RF-5b] Data Sources — Plan B default, Plan A specifics pending Sandia meeting — verified by: AC-6

\paragraph{Software engineering corpus.} We will use the prior work corpus from~\cite{cusati2026fose}, extended with HPC-focused publications from arXiv (cs.DC, cs.SE) and IEEE Xplore, and seeded with Graph-Based RAG~\cite{Edge2024GraphRAG} for the Phase 1 calibration study described in \S\ref{sec:methods}.

\paragraph{HPC codebases and performance data.} We will anchor the HPC evaluation on Sandia-based scientific software systems that together exercise the full performance-portability-to-application stack. The \textbf{initial implementation effort} focuses on \textbf{Pyomo}~\cite{hart2017pyomo}---a Python-based open-source software package originally developed at Sandia National Laboratories for formulating and solving a broad class of optimization problems, including linear, quadratic, nonlinear, mixed-integer (linear, quadratic, nonlinear, and stochastic), generalized disjunctive, differential algebraic equation, and mathematical programming with equilibrium constraint formulations. Pyomo supports symbolic problem definition, concrete instance construction, and integration with standard solvers within a full-featured Python environment; this architecture also underpins \textbf{mpi-sppy}, a parallel stochastic programming package that leverages Pyomo's HPC-ready modeling objects. Pyomo's breadth of problem types, layered abstraction model, and decades of algorithmic design decisions accumulated through active Sandia development (weekly developer coordination meetings; contact: \texttt{wg-pyomo@sandia.gov}) make it a well-instrumented initial corpus for PA-AKG: it encodes exactly the kind of complex, heterogeneous scientific knowledge---numerical method choices, solver compatibility trade-offs, performance characteristics---that must be explicitly represented in the KG for AI-assisted code recommendations to be trustworthy. The broader Sandia evaluation will additionally target \textbf{Kokkos} (performance portability layer), \textbf{Trilinos} (numerical solver framework), and \textbf{LAMMPS} (molecular dynamics application) as cross-domain generalizability cases.

Public data sources include the Extreme-Scale Scientific Software Stack (E4S) ecosystem, the public LAMMPS GitHub history and pair-style benchmark corpus, public Kokkos tutorials and regression tests, the Trilinos open benchmark suite, and the Pyomo GitHub repository history and associated benchmark artifacts. Performance and compiler-trace data will be drawn from SPEC CPU/HPC benchmarks as a public baseline.
\todo{RF-5b (post-Sandia-meeting): after the 2026-04-09 meeting with Co-PI Steinmetz, swap in the specific Sandia Pyomo/Kokkos/Trilinos/LAMMPS corpus subsets, benchmark logs, and validation artifacts accessible under existing Sandia data agreements. Cite the data agreement and name the specific code subsets.}

\paragraph{Pre-trained models.} For local inference and privacy-sensitive prototyping we will use \textbf{Llama 3} (Meta, open-weight). For comparison and high-capability benchmarks we will evaluate \textbf{GPT-4} (OpenAI API) and \textbf{Claude} (Anthropic API). For code-specific tasks, \textbf{CodeLlama} provides a code-pretrained baseline.

\paragraph{Graph storage and retrieval infrastructure.} The PA-AKG substrate will be implemented on \textbf{Neo4j} (primary development) with Amazon Neptune evaluated as a cloud-deployment target. The dense-retrieval vector index will use \textbf{FAISS} for local experimentation and \textbf{pgvector} for the integrated property-graph + vector-index queries required by the canonicalization-auditing pipeline.

\paragraph{Data access.} Virginia Tech has existing access to public corpora (arXiv, IEEE Xplore, GitHub). Sandia National Laboratories, via Co-PI Steinmetz, provides access to the HPC validation corpora identified above under existing Sandia data agreements; the specific data-transfer mechanism and scope are being finalized in coordination with Sandia information-release procedures.
